% Einheit 1 von 3 – Der Quotient (bank/bedingte-wahrscheinlichkeit-und-bayes/e1.jsonl, zusammenbau v0.1)
%% TODO Zeitmarke Einheit 1: Marken-Zeile nicht lesbar
%% TODO Zweigzeile Teil 1: was hier gelernt wird (Satz in Schülersprache aus dem Einheitstitel) – Einheit 1
\einheitenkopf[e1]{Einheit 1 von 3 · Der Quotient}
\zweigzeile{Abitur GK}
\setcounter{aufgabe}{7}

%% TODO Titel: Ich-kann-Satz fehlt in der Bank (Platzhalter: Kettenname) – Nr. 8
%% TODO Anweisung: ein Satz über der ersten Teilaufgabe fehlt in der Bank – Nr. 8
\begin{aufgabe}{Der Quotient}
%% TODO \rechenplatz{n}: Schrittzahl des Grundfalls fehlt in der Bank (nur bei fester Schreibform, 2.3 b)
\begin{teile}
\teil Gesucht ist der Anteil der Mädchen unter den Vereinsmitgliedern, die Fußball spielen. Durch wessen Anteil wird geteilt? \\ \kreuz{alle Mitglieder} \\ \kreuz{die Fußball spielenden Mitglieder} \\ \kreuz{die Mädchen}
\teil In einem Jahrgang lernen 40\,\% der Schüler Spanisch; 16\,\% aller Schüler lernen Spanisch und spielen ein Instrument. Ein Schüler, der Spanisch lernt, wird zufällig ausgewählt. Mit welcher Wahrscheinlichkeit spielt er ein Instrument? \leerfeld
\teil Befragt wurden 200 Personen. A: „trinkt morgens Kaffee“, B: „steht vor 6 Uhr auf“. Die Vierfeldertafel zeigt die Anzahlen. Eine Person, die vor 6 Uhr aufsteht, wird zufällig ausgewählt. Mit welcher Wahrscheinlichkeit trinkt sie morgens Kaffee?

\vierfeldertafel{A}{B}{36,54,90,24,86,110,60,140,200}
\leerfeld
\teil In einem Jahrgang bedeutet A: „trägt eine Brille“, B: „ist ein Mädchen“. Die Vierfeldertafel zeigt die Anzahlen. Ist der Anteil der Brillenträger unter den Mädchen größer als unter den Jungen?

\vierfeldertafel{A}{B}{15,35,50,14,21,35,29,56,85}
\teil A: „hat eine Allergie“, B: „wohnt in der Stadt“. Die Vierfeldertafel zeigt die Anteile. Berechne $P_B(A)$ und $P_A(B)$ und erkläre, warum die Werte verschieden sind.

\vierfeldertafel{A}{B}{0.12,0.18,0.3,0.28,0.42,0.7,0.4,0.6,1}
\end{teile}
\end{aufgabe}

%% TODO Titel: Ich-kann-Satz fehlt in der Bank (Platzhalter: Kettenname) – Nr. 9
%% TODO Anweisung: ein Satz über der ersten Teilaufgabe fehlt in der Bank – Nr. 9
\begin{aufgabe}{Der Quotient – Fehler finden, Begründen, Anwendung}
\begin{teile}
\teil In einem Hotel buchen 35\,\% der Gäste Frühstück, 28\,\% der Gäste sind Geschäftsreisende und 14\,\% aller Gäste sind Geschäftsreisende mit Frühstück. Gesucht ist die Wahrscheinlichkeit, dass ein Geschäftsreisender Frühstück bucht. Ben rechnet $\frac{0{,}14}{0{,}35} = 0{,}4$. Finde den Fehler.
\teil Begründe, warum beim Quotienten für eine bedingte Wahrscheinlichkeit im Nenner der Anteil der Bedingung steht.
\teil An einer Schule kommen 55\,\% der Schüler mit dem Bus; 22\,\% aller Schüler kommen mit dem Bus und sind in diesem Monat mindestens einmal zu spät gekommen; insgesamt kamen 30\,\% mindestens einmal zu spät. Die Schulleitung vermutet, dass Busfahrer häufiger zu spät kommen als die übrigen Schüler. Stimmt das?
\end{teile}
\end{aufgabe}
